\documentclass[12pt,letterpaper]{article}
\usepackage[T1]{fontenc}
\usepackage{mathptmx}
\usepackage{courier}
\usepackage[margin=1in]{geometry}
\usepackage{amsmath,amssymb,amsthm}
\usepackage{enumitem}
\usepackage{booktabs}
\usepackage{array}

\newcommand{\F}{\mathbb{F}}
\newcommand{\Z}{\mathbb{Z}}
\newcommand{\C}{\mathbb{C}}
\newcommand{\Cent}{\operatorname{Cent}}
\newcommand{\Ind}{\operatorname{Ind}}
\newcommand{\tr}{\operatorname{tr}}
\newcommand{\supp}{\operatorname{supp}}
\newcommand{\PGL}{\mathrm{PGL}}
\newcommand{\PSL}{\mathrm{PSL}}
\newcommand{\GL}{\mathrm{GL}}
\newcommand{\SL}{\mathrm{SL}}

\theoremstyle{definition}
\newtheorem*{fact}{Fact}
\newtheorem*{claim}{Claim}

\setlength{\parskip}{4pt}
\setlength{\emergencystretch}{2em}

\title{Referee report on\\
\emph{The representation theory of finite commutation transitive groups}\\
(D.~Joyner, preliminary version dated 8-5-2000, file \texttt{comm\_tr6.tex})}
\author{}
\date{16 November 2026}

\begin{document}
\maketitle

\section*{Summary of the submission}

The note considers finite groups $G$ in which commutation is an
equivalence relation on $G\setminus\{1\}$ (``commutative transitive'',
CT) or on $G\setminus Z(G)$ (``commutation transitive''). Equivalently
(the Lemma of \S3), the centralizer of every non-central element is
abelian. After recalling the classification of solvable and simple CT
groups, it (i) records a Mackey irreducibility criterion, (ii) gives a
formula for the character of $\Ind_C^G\mu$ when $C$ is a centralizer,
(iii) states a corollary on the reducibility of $\Ind_T^G\pi$ for a
torus $T$ in $\PGL(2,\F_q)$, (iv) defines a ``support map''
$\tau\colon L(G)\to\F_2^{G_*}$ and proves (the main theorem in
the source) that the permutation characters $\Ind_A^G 1$, for
centralizers $A$ with $\F_2$-independent supports, are linearly
independent, (v) conjectures that $\dim_{\F_2}\operatorname{Span}
\{\tau(\Cent(g,G))\}\le 4$ for simple CT groups, and (vi) announces a
profinite analogue.

\section*{Overall assessment}

The author is right that there are serious errors. I found:
\begin{itemize}[leftmargin=*,itemsep=1pt]
\item one false example ($\PGL(2,\F_q)$, $q$ odd, is \emph{not} CT);
\item a wrong character formula, off by a factor $|C|$ (and also wrong at
  non-identity central elements);
\item a corollary that refers to a non-existent Proposition, applies a
  lemma outside its hypotheses, is false for $q=2$, and is trivially true
  for $q\ge3$ for a reason unrelated to its hypothesis;
\item an invalid proof of the main theorem. The result is true in a
  corrected form, and the corrected proof is two lines long and needs no
  $\F_2$-linear algebra;
\item a conjecture that, with the definitions as written, is an easy
  theorem. The dimension is exactly $3$ for every $\PSL(2,\F_{2^k})$ and
  is at most $3$ for every finite CT group;
\item a concluding claim (all of $G^*$ from linear combinations of
  monomial representations) that does not follow and is false by a
  dimension count;
\item an empty final section.
\end{itemize}
All of these can be fixed. But once they are fixed, almost nothing new
is left: the surviving statements are short exercises on top of the
classical classification. In its present form the note is not
publishable. Under \emph{Feasibility} below I suggest what a viable
revision might look like.

I checked the group-theoretic claims below by brute force in Python, for
$\GL(2,q)$, $\PGL(2,q)$, $\PSL(2,q)$ with $q\in\{2,3,5,7\}$ and
$\SL(2,2^k)$ with $k\in\{2,3,4\}$. The numbers quoted come from those
runs.

\section{Major mathematical issues}

\subsection*{M1. Example (``$\PGL(2,\F_q)$, $q$ odd, is commutative transitive'') is false}

The counterexample is the author's own footnote to the Remark after
Theorem~2. Let $q$ be odd, $t=\overline{\mathrm{diag}(-1,1)}$,
$w=\overline{\left(\begin{smallmatrix}0&1\\1&0\end{smallmatrix}\right)}$
and $\overline{T}$ the image of the diagonal torus. Then
$w^{-1}tw=\overline{\mathrm{diag}(1,-1)}=t$, so
$\Cent(t,\PGL(2,\F_q))\supseteq\langle\overline T,w\rangle\cong
D_{2(q-1)}$. This dihedral group is non-abelian once $q-1\ge3$. For
$q=3$, $\PGL(2,\F_3)\cong S_4$ and $\Cent((12)(34))\cong D_8$. Since
$Z(\PGL(2,\F_q))=1$, $\PGL(2,\F_q)$ is neither commutative transitive
nor commutation transitive for any odd $q$. (Brute-force check:
$\PGL(2,q)$ fails for $q=3,5,7$, while $\GL(2,q)$ passes; $\PSL(2,5)\cong
A_5$ and $\PSL(2,3)\cong A_4$ pass, $\PSL(2,7)$ fails.)

This is exactly the mechanism the Remark describes: commutation in
$G/Z$ is not commutation in $G$. The Example should be deleted. If an
example of a non-simple, non-solvable group is wanted, the correct
statement is that $\GL(2,k)$ and $\SL(2,k)$ are commutation transitive
for every field $k$. That is Theorem~2(2), whose one-line proof should
be given: a non-scalar $g\in M_2(k)$ is cyclic, so its centralizer in
$M_2(k)$ is the commutative algebra $k[g]$.

The parenthetical facts in the Example are correct but beside the
point:
\[
[\PGL(m,\F_q):\PSL(m,\F_q)]=\gcd(m,q-1),
\]
which equals $2$ when $m=2$ and $q$ is odd.

\subsection*{M2. The induced-character formula (Lemma before the Corollary) is wrong}

Let $C=\Cent(g_0,G)$ with $g_0\notin Z$, and $\mu\in C^*$. The standard
formula is
\[
\chi(g)=\frac{1}{|C|}\sum_{\substack{x\in G\\ x^{-1}gx\in C}}\mu(x^{-1}gx).
\]
\begin{enumerate}[label=(\alph*),itemsep=1pt]
\item \emph{Non-central $g$.} Suppose $a\in\operatorname{Conj}(g,G)\cap
  C$, say $g=sas^{-1}$. Then $a\notin Z$, so $\Cent(a,G)=C$ by
  the centralizer lemma of \S2. For $u=s^{-1}x$ we have $u^{-1}au\in C$ if
  and only if $u^{-1}Cu=\Cent(u^{-1}au,G)=C$, that is, $u\in N_G(C)$.
  Since $C$ is abelian,
  \[
  \tr\pi(g)=\frac1{|C|}\sum_{n\in N_G(C)}\mu(n^{-1}an)
  =\sum_{w\in W}\mu(w^{-1}aw),\qquad W:=N_G(C)/C.
  \]
  The paper omits the factor $|C|^{-1}$. Concretely, take $G=S_3$,
  $C=A_3$, $\mu$ nontrivial. At a $3$-cycle the true value is $-1$; the
  paper's formula gives $-3$.
\item \emph{Central $g\ne1$}, which occurs when $Z\ne1$ (for example in
  $\GL(2,q)$). Here $\tr\pi(g)=\tfrac{|G|}{|C|}\mu(g)$, the same shape
  as at $g=1$. The paper's first line instead gives $|N_G(C)|\mu(g)$.
\item The argument in (a) also shows that the elements of
  $\operatorname{Conj}(g,G)\cap C$ form one $W$-orbit. This is exactly
  what makes ``$a=a_g$'' well defined, and the paper should say so.
\end{enumerate}
\emph{Fix:} replace the Lemma with the three-case formula above, using
$W=N_G(C)/C$. The proof is the five lines in (a). Note that the
``Proof'' cites \cite[Ch.~7]{Se}, but the result needs only the
definition of the induced character.

\subsection*{M3. The Corollary on $\PGL(2,\F_q)$}

\begin{enumerate}[label=(\alph*),itemsep=1pt]
\item It is labelled ``(of Proposition)'', but there is no Proposition
  in the paper.
\item Its proof uses the Frobenius-formula Lemma, which assumes $G$ is
  commutation transitive. By M1, $\PGL(2,\F_q)$ is not when $q$ is odd.
  The specific step that fails is $\Cent(a,G)=C$ for $a\in C\setminus
  Z$: at the involution $t\in\overline T$ one has
  $\Cent(t)=N(\overline T)\ne\overline T$.
\item \emph{The statement is false for $q=2$.} $\PGL(2,\F_2)\cong S_3$,
  the nonsplit torus is $T=A_3$, $N_G(T)=S_3$, and for nontrivial $\pi$
  we have $\pi^h=\pi^{-1}\neq\pi$ for all $h\in H\setminus T$. Yet
  $\Ind_{A_3}^{S_3}\pi$ is the irreducible $2$-dimensional
  representation.
\item For $q\ge3$ the conclusion holds for every $\pi$, without the
  hypothesis, because of degrees: $[G:T]=q(q\pm1)>q+1$, the largest
  irreducible degree. In fact Mackey's formula gives exact values. Every
  $T$--$T$ double coset outside $N_G(T)$ has trivial intersection
  $T\cap T^s$, and $[N_G(T):T]=2$. Writing $\delta=1$ if $\pi=\pi^w$
  and $0$ otherwise,
  \[
  \langle\Ind_T^G\pi,\Ind_T^G\pi\rangle=
  \begin{cases}
  q+3+\delta, & T\text{ split},\\
  q-1+\delta, & T\text{ nonsplit}.
  \end{cases}
  \]
  I confirmed this numerically for split $T$ with $q=3,5,7$ and for
  nonsplit $T$ with $q=2,3$.
\item The proof itself is not an argument. It mixes $\pi$ and $\mu$.
  The ``cancellation'' it asserts is not justified, since the double sum
  is not a character sum over a group in any evident way. And the
  assertion that the $g=1$ terms contribute more than $1$ fails exactly
  at the counterexample: there $|G|/|T|^2=6/9$.
\end{enumerate}
\emph{Fix:} replace the Corollary and its proof with the Mackey
computation in (d). This is routine, standard material (compare
\cite{B} or any treatment of $\GL_2(\F_q)$). It belongs in a
``non-example'' remark, since $\PGL(2,\F_q)$ is not in the class being
studied.

\subsection*{M4. The main theorem: invalid proof; the correct version is easier}

\emph{The proof.} ``The functions are supported on pairwise distinct
sets, so no linear combination can vanish'' is false as a principle.
Take three classes $\alpha,\beta,\gamma$ and $f_1=1_\alpha+1_\beta$,
$f_2=1_\alpha+1_\gamma$, $f_3=2\cdot1_\alpha+1_\beta+1_\gamma$. Their
supports $\{\alpha,\beta\}$, $\{\alpha,\gamma\}$,
$\{\alpha,\beta,\gamma\}$ are distinct, and even independent over $\F_2$
(the $0/1$ matrix has determinant $-1$). But $f_3=f_1+f_2$. So
$\F_2$-independence of the $\tau(A)$ neither implies nor is implied by
$\C$-independence of the characters. For the latter direction: in
$A_5\cong\PSL(2,\F_4)$ the four centralizer classes $V_4,C_3,C_5,G$
satisfy $\tau(G)=\tau(V_4)+\tau(C_3)+\tau(C_5)$ over $\F_2$, yet the
four permutation characters $1_A^G$ are $\C$-linearly independent
(rank~$4$, computed).

\emph{The statement.} Three problems:
\begin{enumerate}[label=(\roman*),itemsep=1pt]
\item The set $\mathcal C$ is indexed ambiguously. Conjugate centralizers
  have equal $\tau$, and equal induced characters, so as written either
  $\mathcal C$ contains repeats (and the conclusion is false) or the
  indexing must be one centralizer per conjugacy class.
\item $\chi_A$ is defined as the trivial character, while the theorem
  writes ``$\chi=\chi_\gamma\in A^*$''. It is unclear whether arbitrary
  characters are intended.
\item The symbol $\chi$ already denotes a central character in the
  Mackey lemma.
\end{enumerate}

\emph{A correct statement with a short proof.} Let $G$ be finite and
commutation transitive, and let $A_1,\dots,A_m$ be representatives of
the conjugacy classes of centralizers of non-central elements. A class
$\gamma\not\subset Z$ meets $A_i$ if and only if $A_i$ is conjugate to
$\Cent(g,G)$ for $g\in\gamma$. So the sets $D_i:=\supp(A_i)\setminus
\{\text{central classes}\}$ are non-empty and pairwise disjoint. On
$G\setminus Z$, the function $1_{A_i}^G$ is supported on $D_i$, where it
is positive. Hence $1_{A_1}^G,\dots,1_{A_m}^G$ (and also the trivial
character, by evaluating at $1$ or at a central element) are linearly
independent. The same argument shows that the $\tau(A_i)$ are
\emph{automatically} $\F_2$-independent. The $\F_2$ machinery can be
dropped, and $\mathcal C$ is simply ``one centralizer from each
conjugacy class''.

The catch is that $m$ is tiny. One has $m=2$ for a Frobenius group with
abelian kernel and cyclic complement, $m=3$ for $\PSL(2,\F_{2^k})$
(unipotent $U$, split torus, nonsplit torus), and $m=3$ for $\GL(2,q)$
($ZU$ and the two tori). So the theorem yields at most four independent
characters.

\subsection*{M5. The Conjecture is an easy theorem, with $3$ in place of $4$}

Let $G=\PSL(2,\F_q)$ with $q=2^k\ge4$. Every non-identity element is
exactly one of: unipotent, with centralizer a conjugate of $U$ (order
$q$, since $-1=1$); split regular semisimple, with centralizer a
conjugate of $T$ (order $q-1$); or nonsplit regular semisimple, with
centralizer a conjugate of $T'$ (order $q+1$). So
$\{\tau(\Cent(g,G))\}=\{\tau(G),\tau(U),\tau(T),\tau(T')\}$. Over
$\F_2$, the identity class occurs three times in
$\tau(U)+\tau(T)+\tau(T')$ and every other class once, giving
\[
\tau(U)+\tau(T)+\tau(T')=\tau(G).
\]
The first three are independent because their non-identity parts are
disjoint and non-empty. Hence the dimension is exactly~$3$ for all~$k$.
This agrees with brute force:

\begin{center}
\begin{tabular}{lccc}
\toprule
$G$ & $|G_*|$ & distinct $\tau$-vectors & $\F_2$-rank\\
\midrule
$\PSL(2,\F_4)$ & 5 & 4 & 3\\
$\PSL(2,\F_8)$ & 9 & 4 & 3\\
$\PSL(2,\F_{16})$ & 17 & 4 & 3\\
\bottomrule
\end{tabular}
\end{center}

The ``$\le4$'' observed in the GAP runs is probably the number of
\emph{distinct} vectors rather than the rank; the author should recheck
the code. More generally, the M4 argument gives $\dim=m$ when $m$ is odd
and $m+1$ when $m$ is even. By the classification (abelian, Frobenius,
$\PSL(2,\F_{2^k})$) this is at most $3$ for every finite CT group, not
only the simple ones. The conjecture should become a proposition with
this proof, and it is not deep.

\subsection*{M6. ``All representations in $G^*$ can be obtained from linear combinations of monomial representations''}

This does not follow from the main theorem, which asserts independence,
not spanning. It is also false for the characters considered. They span
a space of dimension at most $m+1\le4$, while $|G^*|=|G_*|=q+1$ for
$\PSL(2,\F_q)$, which is already $5$ for $A_5$. The true statement in
this direction is Brauer's induction theorem: every character is a
$\Z$-combination of characters induced from linear characters of
elementary subgroups. Brauer's subgroups are generally not centralizers.
The attribution to Gelfand--Graev is inappropriate. Gelfand--Graev
representations are induced from a non-degenerate character of a
maximal unipotent subgroup and are multiplicity-free; that is a
different phenomenon. The abstract's promise to determine ``which
irreducible characters \dots\ are generated by certain monomial
representations'' is not carried out anywhere.

\subsection*{M7. Mackey's criterion and the profinite section}

\begin{enumerate}[label=(\alph*),itemsep=1pt]
\item In the $p$-adic setting the relevant induction is compact
  induction $c\text{-}\Ind_H^G$, and the conclusion is usually stated as
  ``irreducible and supercuspidal''. Condition~(2), on $|\chi|$
  extending to $G$, is non-standard as stated, and its role should be
  explained or a reference given. The conventions for $H^g$ and $\pi^g$
  are never defined, and with the paper's $y^{-1}xy$ convention the
  direction matters. ``Mauntner'' should be ``Mautner''.
\item For finite groups the criterion is an \emph{if and only if}
  \cite[\S7.4, Prop.~23]{Se}. That is what justifies concluding
  reducibility in Remark~(b), and the paper should say so.
\item Remark~(b): for $G=\GL(2,\F_q)$ the claim $T\cap T^g\in\{T,Z\}$
  is correct, because non-central torus elements are regular. Note that
  $T\cap T^g=T$ holds for $g\in N_G(T)\setminus T$, not only $g\in T$.
  For $\PSL(2,\F_q)$, ``$Z\cong\F_q^\times$'' is wrong, since the center
  is trivial. For odd $q\ge7$, $\PSL(2,\F_q)$ is not even in the class
  under study. The footnote ``Disjoint groups do not have disjoint
  representations!'' should be replaced by the plain statement that
  $\pi$ and $\pi^g$ have the same restriction to $Z$.
\item \S5 states, without formulation or proof, that ``there is an
  analog'' for profinite CT groups. Either supply a theorem with a proof
  (for example via continuous characters factoring through finite
  quotients, being careful that centralizers in a profinite group are
  closed but need not be open), or delete the section and the sentence
  in the abstract.
\end{enumerate}

\section{Minor mathematical issues}

\begin{enumerate}[leftmargin=*,itemsep=2pt]
\item \emph{After the centralizer lemma of \S2:} ``any two centralizers \dots\
  are either disjoint groups or equal''. Subgroups are never disjoint.
  The correct statement is $C_1\cap C_2=Z$ or $C_1=C_2$. This is also
  what is needed for the Frobenius-complement argument.
\item \emph{The bijection between $\operatorname{Conj}(g,G)$ and the
  cosets of $\Cent(g,G)$:} with the convention $x^{-1}gx$, the map
  $x^{-1}gx\mapsto x\Cent(g,G)$ is not well defined. The correct
  bijection is with right cosets, $\Cent(g,G)x$. Alternatively, use
  $xgx^{-1}\leftrightarrow x\Cent(g,G)$. (Checked in $S_3$.)
\item \emph{Remark after Theorem~2:} the relation $\sim$ transported to
  $(G-Z)/Z$ \emph{is} transitive, since it is equivalent to $\sim$ on
  $G-Z$. What can fail is that commutation in $G/Z$ coincides with this
  $\sim$, and transitivity of commutation in $G/Z$ (see M1). The
  sentence ``so transitivity can fail'' should say which relation is
  meant. The notation $x^y$ is used before it is defined.
\item \emph{The Frobenius-subgroup paragraph in \S2.}
  \begin{itemize}[itemsep=0pt]
  \item The condition $C\cap C^x=1$ for $x\notin N_G(C)$ needs $Z=1$.
    In general one should pass to $C/Z\le G/Z$, where it does hold.
  \item ``A normal subgroup of size at least $2$'' is vacuous, since $G$
    itself is normal. The intended statement is a proper nontrivial
    normal subgroup: the Frobenius kernel.
  \item Please give the precise location of Frobenius' theorem in
    \cite{Se}.
  \end{itemize}
\item \emph{Theorem~1.}
  \begin{itemize}[itemsep=0pt]
  \item ``Cyclic'' should be ``cyclic of prime order''.
  \item ``A cyclic group $H$ of fixed-point-free automorphisms'' should
    require $H\ne1$ and every non-identity element of $H$ to act without
    nontrivial fixed points; that is, $G$ is a Frobenius group with
    abelian kernel and cyclic complement.
  \item The paper treats a ``general case'', so it should also state the
    nonsolvable case: as I understand the literature (Suzuki;
    Brauer--Suzuki--Wall; Wu), a finite nonsolvable CT group is
    isomorphic to $\PSL(2,\F_{2^k})$, $k\ge2$. The author should confirm
    the exact statement in \cite{Wu}. With it, \S4.2 reduces to two
    explicit families, which would make the paper much more concrete.
  \end{itemize}
\item \emph{\S4.1:} ``not cyclic'' should be ``non-abelian''. The
  little-groups method of \cite[\S8.2]{Se} needs a semidirect product
  with an \emph{abelian} normal factor. That is available for solvable
  CT groups, but not for solvable commutation transitive groups in
  general. For example, $\GL(2,\F_3)$ is solvable and commutation
  transitive, but its only normal abelian subgroups are $1$ and $Z$,
  neither of which has a complement.
\item \emph{Terminology.} The two notions are used interchangeably in
  several places:
  \begin{itemize}[itemsep=0pt]
  \item after the centralizer lemma of \S2;
  \item the Remark in \S3;
  \item \S4.1;
  \item the main theorem, stated for commutation transitive groups,
    versus the Conjecture, stated for CT groups.
  \end{itemize}
  Fix one name for each notion and use it consistently. Groups in which
  every non-central centralizer is abelian are also widely called
  \emph{AC-groups} (or CA-groups). The author should check the classical
  work on them; I believe R.~Schmidt (1970) classified the finite ones,
  but this should be verified. The graph $\Gamma$ of \S2 is the
  \emph{commuting graph}, which has a substantial literature.
\item \emph{Lemma in \S2(b):} state the bijection explicitly, $K\mapsto
  K\cup Z$ with inverse $C\mapsto C\setminus Z$. Also, $\Cent$ is used
  before it is defined.
\end{enumerate}

\section{Bibliography}

\begin{itemize}[leftmargin=*,itemsep=1pt]
\item \cite{S}: I believe the title of Suzuki's 1957 Proc.\ AMS paper is
  ``\dots\ simple groups of \emph{odd} order'', not ``even order''.
  Please check.
\item ``Wall'' is credited in the introduction but not cited. Presumably
  the intended reference is R.~Brauer, M.~Suzuki, G.~E.~Wall, \emph{A
  characterization of the one-dimensional unimodular projective groups
  over finite fields}, Illinois J.\ Math.\ \textbf{2} (1958).
\item \cite{J}: ``Archiv der Math.\ 273(1999)'' cannot be right, since
  the 1999 volumes are numbered in the 70s. Check the volume and pages.
\item \cite{B} is never cited. \cite{FGRS} (``to appear'') and \cite{F}
  (``preprint 1983'') should be updated.
\item Typos: ``Klienian'' should be ``Kleinian''; ``N, Iwahori'' should
  be ``N.~Iwahori''. In \cite{LR}, check whether the title reads
  ``commutation-transitive''.
\end{itemize}

\section{Readability and presentation}

\begin{itemize}[leftmargin=*,itemsep=1pt]
\item \emph{Typos:} ``autmorphisms'', ``commutivity'', ``whcih'',
  ``methif'', ``power set offfff'', ``this papers'', ``Mauntner''.
  ``Commutative transitive'' and ``commutation transitive'' are also
  swapped in places (see minor item~7).
\item \emph{Undefined notation:} $H^g$, $\pi^g$, $x^y$, $a_g$, $L(G)$
  (``subgroup lattice'' is mentioned only after first use), and
  $\chi_\gamma$. The symbols $\pi$/$\mu$ and $\chi$ are overloaded.
\item \emph{Structure.} The Lemma in \S3 is the working definition of
  the paper and should come first, in \S1. The Remarks after Theorems~1
  and~2 read as notes to self. ``Proof: Trivial'' and
  ``immediate'' should be replaced by the one or two lines that are
  actually needed.
\item \emph{Abstract:} it should describe only what the paper proves.
  At present the second and third sentences do not correspond to proved
  results.
\item \emph{\TeX:} \verb|\Bbb| is obsolete (use \verb|\mathbb|), the
  class option \verb|amsfonts| is unused, there is a stray ``b'' after
  \verb|\end{document}|, and there is a missing parenthesis in
  ``$\mathbb{F}_q$-valued points of a maximal torus of $PGL(2)$''. The
  file compiles, apart from these warnings.
\end{itemize}

\section{Feasibility of repair}

\begin{center}
\renewcommand{\arraystretch}{1.15}
\begin{tabular}{>{\raggedright}p{2.7cm}>{\raggedright}p{6.8cm}p{3.6cm}}
\toprule
Item & Repair & Effort\\
\midrule
M1 Example & Delete; replace with $\GL(2,k)$, $\SL(2,k)$ (one-line proof). & Trivial\\
M2 Formula & Replace with the corrected formula with $W=N_G(C)/C$ and central case. & Easy\\
M3 Corollary & Replace with the Mackey count $q+3+\delta$, $q-1+\delta$; move to a remark on non-examples. & Easy\\
M4 Main theorem & Restate (one centralizer per conjugacy class); two-line diagonal proof; drop the $\F_2$ map. & Easy, but the result becomes modest\\
M5 Conjecture & Becomes a proposition: $\dim=3$ for $\PSL(2,\F_{2^k})$, $\le3$ for all CT groups. & Easy\\
M6 Spanning claim & Delete, or replace by a discussion of Brauer induction. & Easy (delete); genuine work for a new result\\
M7 Mackey / profinite & Fix the statement and conventions; either prove a profinite theorem or delete \S5. & Moderate to hard if kept\\
Minor items, bibliography, presentation & As listed. & Easy\\
\bottomrule
\end{tabular}
\end{center}

\paragraph{What would remain.} After these repairs, the note consists
of the classification (quoted), the observation that the commuting
graph on $G\setminus Z$ is a disjoint union of cliques, the corrected
induced-character formula, and an easy independence result for at most
four permutation characters. Every finite CT group is abelian, a
Frobenius group $N\rtimes H$ with $N$ abelian and $H$ cyclic, or
$\PSL(2,\F_{2^k})$. All of these have classically known character
tables, so a paper whose aim is ``the representation theory of finite CT
groups'' needs a sharper question to be new.

One natural option is to decompose $\Ind_C^G\mu$ explicitly, for every
centralizer $C$ and every $\mu\in C^*$, in the two nonabelian families.
Then determine which irreducibles occur and with what multiplicity, and
which $\Z$-combinations of these monomial characters give all of $G^*$.
For $\PSL(2,\F_{2^k})$ this is a pleasant exercise with the known
character table. It would give real content to the abstract's second
sentence. A second option is a commutation transitive analogue, with
$Z\ne1$, for $\GL(2,\F_q)$ and $\SL(2,\F_q)$. That is where the
$p$-adic motivation (Mackey/Mautner) actually lives.

\paragraph{Recommendation.} Not acceptable in its present form. The
errors are real but local and quickly fixed. The main obstacle is
content rather than correctness: the corrected main results are short
exercises. I would encourage either (a) recasting the material as a
short expository note, perhaps a blog post or a GAP/Sage worksheet,
with the corrections above, or (b) developing one of the two
directions in the preceding paragraph into a genuine paper.

\begin{thebibliography}{FGRS}
\bibitem[B]{B} A.~Borel et al., \emph{Seminar on algebraic groups and related finite groups}, LNM 131, Springer.
\bibitem[F]{F} B.~Fine, preprint 1983 (as cited in the submission).
\bibitem[FGRS]{FGRS} B.~Fine, A.~Gaglione, G.~Rosenberger, D.~Spellman, \emph{The commutative transitive kernel} (as cited).
\bibitem[J1]{J} D.~Joyner, Archiv der Math.\ (as cited; volume to be checked).
\bibitem[LR]{LR} F.~Levin, G.~Rosenberger, in \emph{Groups -- St.\ Andrews 1985}, LMS LN 121.
\bibitem[S]{S} M.~Suzuki, Proc.\ AMS \textbf{8} (1957), 686--695.
\bibitem[Se]{Se} J.-P.~Serre, \emph{Linear representations of finite groups}, GTM 42, Springer, 1977.
\bibitem[Wu]{Wu} Y.-F.~Wu, \emph{Groups in which commutativity is a transitive relation}, J.\ Algebra \textbf{207} (1998), 165--181.
\end{thebibliography}

\noindent\emph{Note on labels.} References such as ``the centralizer lemma of \S2''
and ``the main theorem'' are to the labels in the submitted source
(the centralizer lemma in \S2 and the main theorem in \S4.2).

\end{document}
